% we include the glossary here (frontmatter is included with \input, so this command is as if it was in main.tex)
%All acronyms must be written in this file.
\newacronym{RTS}{RTS}{Real-Time System}
\newacronym{GPOS}{GPOS}{General Purpose Operating System}
\newacronym{RTOS}{RTOS}{Real-Time Operating System}
\newacronym{PGF}{PGF}{Portable Graphics Format}

% Common acronyms used in the thesis
\newacronym{CICD}{CI/CD}{Continuous Integration / Continuous Delivery}
\newacronym{MTTD}{MTTD}{Mean Time To Detection}
\newacronym{MTTR}{MTTR}{Mean Time To Recovery}
\newacronym{PRISMA}{PRISMA}{Preferred Reporting Items for Systematic Reviews and Meta-Analyses}
\newacronym{IP}{IP}{Internet Protocol}
\newacronym{DORA}{DORA}{DevOps Research and Assessment}
\newacronym{GDPR}{GDPR}{General Data Protection Regulation}
\newacronym{eBPF}{eBPF}{extended Berkeley Packet Filter}


\frontmatter % Use roman page numbering style (i, ii, iii, iv...) for the pre-content pages

\pagestyle{plain} % Default to the plain heading style until the thesis style is called for the body content

%----------------------------------------------------------------------------------------
%	TITLE PAGE
%----------------------------------------------------------------------------------------

\maketitlepage


%----------------------------------------------------------------------------------------
%	STATEMENT of INTEGRITY
%----------------------------------------------------------------------------------------
\integritystatement

%----------------------------------------------------------------------------------------
%	DEDICATION  (optional)
%----------------------------------------------------------------------------------------
%
%\dedicatory{For/Dedicated to/To my\ldots}
% \begin{dedicatory}
% 	The dedicatory is optional. Below is an example of a humorous dedication.

% 	"To my wife Marganit and my children Ella Rose and Daniel Adam without whom this book would have been completed two years earlier." in "An Introduction To Algebraic Topology" by Joseph J. Rotman.
% \end{dedicatory}

%----------------------------------------------------------------------------------------
%	ABSTRACT PAGE
%----------------------------------------------------------------------------------------

\begin{abstract}

	% here you put the abstract in the main language of the work.
Containerized \ac{CICD} platforms deployed on\-premise with non\-privileged access present critical observability challenges. Organizations managing infrastructure without host administrative privileges cannot deploy standard monitoring solutions that assume privileged container execution or cloud\-managed telemetry. This research investigates container\-first observability approaches for heterogeneous environments where conventional instrumentation is unavailable. Through systematic literature review following \ac{PRISMA} methodology, the work identifies a gap between cloud\-native observability assumptions and operational realities in privilege\-restricted deployments. The research designs and evaluates an integrated observability framework collecting metrics, logs and distributed traces from non\-privileged containers across mixed Linux and Windows hosts. The solution employs sidecar exporters, agent collectors and runtime APIs to instrument \ac{CICD} services without modifying container images or requiring host\-level access. Empirical evaluation measures \ac{MTTD}, \ac{MTTR}, telemetry pipeline throughput and instrumentation overhead against a defined baseline using controlled fault injection. The primary contributions include a reproducible reference architecture, prototype implementation with deployment automation, and practical guidance for constrained environments. The work provides evidence\-based assessment of whether container\-first instrumentation can measurably improve fault detection and diagnosis while maintaining resource overhead below operationally acceptable thresholds in privilege\-restricted \ac{CICD} platforms.
\end{abstract}


%----------------------------------------------------------------------------------------
%	ACKNOWLEDGEMENTS (optional)
%----------------------------------------------------------------------------------------

% \begin{acknowledgements}

% 	The optional Acknowledgment goes here\ldots Below is an example of a humorous acknowledgment.

% 	"I'd also like to thank the Van Allen belts for protecting us from the harmful solar wind, and the earth for being just the right distance from the sun for being conducive to life, and for the ability for water atoms to clump so efficiently, for pretty much the same reason. Finally, I'd like to thank every single one of my forebears for surviving long enough in this hostile world to procreate. Without any one of you, this book would not have been possible." in "The Woman Who Died a Lot" by Jasper Fforde.
% \end{acknowledgements}

%----------------------------------------------------------------------------------------
%	LIST OF CONTENTS/FIGURES/TABLES PAGES
%----------------------------------------------------------------------------------------

\tableofcontents % Prints the main table of contents

\listoffigures % Prints the list of figures

\listoftables % Prints the list of tables

% Keep English name by default; remove conditional Portuguese redefinitions
% \renewcommand{\listalgorithmname}{List of Algorithms}
% \listofalgorithms % Prints the list of algorithms
% \addchaptertocentry{\listalgorithmname}


% \renewcommand{\lstlistlistingname}{List of Source Code}
% \renewcommand{\lstlistlistingname}{List of Source Code}
% \lstlistoflistings % Prints the list of listings (programming language source code)
% \addchaptertocentry{\lstlistlistingname}


%----------------------------------------------------------------------------------------
%	ABBREVIATIONS
%----------------------------------------------------------------------------------------
%\begin{abbreviations}{ll} % Include a list of abbreviations (a table of two columns)
%%\textbf{LAH} & \textbf{L}ist \textbf{A}bbreviations \textbf{H}ere\\
%%\textbf{WSF} & \textbf{W}hat (it) \textbf{S}tands \textbf{F}or\\
%\end{abbreviations}

%----------------------------------------------------------------------------------------
%	SYMBOLS
%----------------------------------------------------------------------------------------




%----------------------------------------------------------------------------------------
%	ACRONYMS
%----------------------------------------------------------------------------------------

\newcommand{\listacronymname}{List of Acronyms}

%Use GLS
\glsresetall
\printglossary[title=\listacronymname,type=\acronymtype,style=long]

%----------------------------------------------------------------------------------------
%	DONE
%----------------------------------------------------------------------------------------

\mainmatter % Begin numeric (1,2,3...) page numbering
\pagestyle{thesis} % Return the page headers back to the "thesis" style
